\subsection{Replies to comments received from Hayato-san - May 19, 2013}


\begin{enumerate}
%........................................................................................................
\vspace{5 mm}
\item
{\em
 I think the proposed definition of 90\% C.L. for $sin^2 2\theta_{23}$ is not appropriate
 because the $\chi^2$ distribution does not have bottom point in the region as shown in the collaboration meeting.
 If you want to use the constant $\chi^2$ method, $\Delta \chi^2$ should have bottom ( minimum ) 
 or change the $\Delta \chi^2$ value from usual number. 
 ( For MINOS case, they have best fit point in the physics region and it may be fine but not for our case. )

 As you find the plot from Wongjirad-san, 2-flavor assumption 90\% C.L. allowing unphysical region
 is almost same as the one you are trying to define as the 90\% only for physically allowed region.
 I think this implies that the proposed method is underestimating the confidence level.

 Therefore, if you believe you can not perform Feldman-Cousins method for this determination,
 it is safe to take the edge of the 2D contour even if the definition is different from the one used by the papers from the other experiment.
}\\
\\
The value of $\chi^{2}$ at  $\sin^{2} 2 \theta_{23}$ = 1 is the minimum since $\chi^{2}$ is not defined for 
$\sin^{2} 2 \theta_{23}$ $>$ 1. This is because, in the 3-flavour analysis, it is not possible to calculate the 
oscillation probabilities for $\sin^{2} 2 \theta_{23}$ $>$ 1 due to the presence of a $\sqrt{(1 - \sin^{2} 2 \theta_{23})}$ term. 
%In a 2-flavour analysis, $\sin^{2} 2 \theta_{23}$ is the amplitude of the oscillation; if it is considered to be a 
%number rather than the square of a trigonometric function, it can take values $>$ 1. However this does not make physical sense 
%since it would mean that there exists an energy at which more than 100\% of neutrinos oscillate. Of course a statistical 
%fluctuation in the data can result in the best-fit value of $\sin^{2} 2 \theta_{23}$ in a 2-flavour analysis being $>$ 1.0. 
%But this is giving information about that particular dataset only; it does not mean that the true value of $\sin^{2} 2 \theta_{23}$ 
%is $>$ 1 since this is physically impossible.

We would also like to point out that, if the minimum value of $\chi^{2}$ in the 1D fit of $\sin^{2} 2 \theta_{23}$ were to be 
lowered, the dashed lines representing the 68\% and 90\% confidence limits would also have lower values. As a result, the lower 
bound of the 90\% confidence interval in $\sin^{2} 2 \theta_{23}$ would be \emph{increased} from 0.95, and that confidence interval 
would actually be narrower. 
 
It is not possible to use the edge of the 2D contour for a 1D limit as they mean different things. The 2D 90\% contour is constructed 
so that, if the experiment were repeated many times, the true unknown values of \emph{both $\sin^{2} 2 \theta_{23}$ and $|\Delta m^{2}_{32}|$} 
would be contained in 90\% of those intervals. By contrast, a 1D interval in $\sin^{2} 2 \theta_{23}$ should be constructed so that 
\emph{$\sin^{2} 2 \theta_{23}$ only} is contained in 90\% of all confidence intervals, irrespective of the true value of $|\Delta m^{2}_{32}|$. 
Furthermore, the 2D contour was made using critical values of $\Delta \chi^{2}$ obtained by fitting many toy MC experiments with both 
$\sin^{2} 2 \theta_{23}$ and $|\Delta m^{2}_{32}|$ and all the systematics. These critical values are higher than those that would be obtained 
from fitting one oscillation parameter and all the systematics to the same toy datasets, and cannot therefore be used for a 1D confidence interval.

\end{enumerate}
